\section{Conclusion}
\label{sec:conclusion}

\subsection{Summary}

We have compared all eight redistricting ensemble algorithms implemented
in \texttt{redist-ensemble} (G.6--G.13) on three states spanning
a representative range of district counts: NC ($k = 14$), WI ($k = 8$),
and TX ($k = 38$).
The main findings are:

\begin{enumerate}

  \item \textbf{No algorithm dominates all metrics.}
        Short-Burst achieves the lowest edge-cut (most compact final plan)
        but samples from a concentrated near-compact distribution.
        SMC achieves the highest mean Polsby-Popper and an exactly
        calibrated target distribution but does not optimise compactness.
        Multi-scale achieves the fastest mixing for TX but adds overhead
        that is unnecessary for small $k$.

  \item \textbf{The compactness--distribution trade-off is fundamental.}
        Algorithms that optimise for compactness (\sburst{}, \sburstf{})
        necessarily sample from a biased distribution; algorithms that
        target uniform or calibrated distributions (\smc{}, \fr{}, \msp{})
        do not produce the most compact plans.
        ShortBurstForest occupies a middle position: better compactness
        than pure \fr{}, better distributional properties than pure
        \sburst{}.

  \item \textbf{Scale is the primary mixing determinant for large $k$.}
        For TX ($k = 38$), Multi-scale reduces lag-100 autocorrelation by
        27\% relative to Merge-Split at comparable runtime, due to
        coarse-level moves that restructure entire county-district
        assignments in a single step.
        For WI ($k = 8$), Multi-scale provides no meaningful mixing
        benefit.

  \item \textbf{VRA compliance and local sensitivity are orthogonal to
        the compactness--distribution axis.}
        VRA-Aware ReCom is the only algorithm that enforces VRA constraints
        during sampling; no other algorithm can substitute for it when
        minority-majority district preservation is required.
        Flip is the only algorithm designed for local boundary sensitivity
        analysis; its high autocorrelation is a feature, not a defect.

  \item \textbf{The four-question decision tree covers all eight algorithms.}
        Four binary questions---VRA compliance, $k \geq 20$, calibrated
        distribution, compactness optimisation---suffice to identify the
        appropriate algorithm for any redistricting task.

\end{enumerate}

\subsection{The Three-Layer Compositor}

The eight algorithms compared in this paper are all implemented as
\texttt{--search} modes in the \texttt{redist} CLI.
The three-layer compositor (\texttt{--structure},
\texttt{--weights-override}, \texttt{--search}) allows practitioners to
combine any structure algorithm (standard bisect, ApportionRegions,
GeoSection, etc.) with any edge-weight signal (geographic, county, VRA,
proportional) and any search algorithm.
This means that the decision tree in Section~\ref{sec:decision} selects
the search layer only; the structure and weight layers are configured
independently based on the specific redistricting task.

For example, a Texas litigation case might use:
\begin{itemize}
  \item \texttt{--structure prime-factor}: ApportionRegions bisection for
        a legally defensible initial plan.
  \item \texttt{--weights county}: county-respecting edge weights to
        encourage county-whole district assignments.
  \item \texttt{--search multiscale}: Multi-scale MCMC for fast mixing
        across the large-$k$ plan space.
\end{itemize}

This combination is not available in any prior academic comparison paper;
the three-layer compositor is a contribution of the \texttt{redist}
platform itself.

\subsection{Future Work}

\paragraph{CVD and BFS Growth.}
Two geometric algorithms---CVD-based redistricting (B.22) and BFS Growth
(B.23)---are in development.
When implemented, they will be added to this comparison.
Both algorithms target geometric compactness directly (rather than
via Markov chain search), which may outperform \sburst{} on \pp{} while
maintaining a broader distributional footprint.

\paragraph{Multi-run confidence intervals.}
The empirical results in this paper are single-run (seed 42).
A follow-up study with $R = 20$ independent runs per algorithm--state pair
would provide confidence intervals on all metrics and enable formal
statistical comparison.
The decision tree would be updated if any ordering reversals appear under
multi-run averages.

\paragraph{Adaptive parameter selection.}
Each algorithm has internal parameters (burst length, $\alpha$, VRA boost)
that affect performance.
Adaptive selection of these parameters---based on real-time monitoring of
acceptance rates, autocorrelation, or compactness trajectory---is the
subject of B.21 (Adaptive Multi-scale) and is a natural extension for
all other algorithms.

\subsection{Implementation Reference}

All eight algorithms are implemented in:
\begin{itemize}
  \item \texttt{crates/redist-ensemble/src/} --- chain implementations
        (\texttt{shortburst.rs}, \texttt{flip.rs}, \texttt{forest\_recom.rs},
        \texttt{merge\_split.rs}, \texttt{multiscale.rs},
        \texttt{shortburst\_forest.rs}, \texttt{vra\_recom.rs})
  \item \texttt{crates/redist-ensemble/src/smc.rs} --- SMC implementation
  \item \texttt{crates/redist-cli/src/compositor/} --- CLI integration and
        three-layer compositor
\end{itemize}

The decision framework in Table~\ref{tab:decision} is the recommended
reference for algorithm selection in the \texttt{redist} platform.
All runs in this paper are reproducible from the base seed 42 and the
configurations specified in Section~\ref{sec:setup}.
